\chapter{不确定信道与长期运行}
\label{chap:stochastic}
\chaptertakeaway{区分估计误差、随机变化和排队状态；分别用最坏情形、概率约束、样本平均和在线决策表达它们。可靠性概率、算法置信度和长期平均性能是不同的量。}

\section{先决定需要保证什么}
信道状态信息（CSI）不完整时，优化对象取决于设计者知道什么、允许什么。若只知道误差位于一个确定集合内，最坏情形设计要求集合内每个信道都满足约束；第4章的范数不确定性属于这种情形。若有可信的概率模型，可以优化平均速率，或允许一个给定的小概率失效。平均性能好，并不能自动保证每个用户的失效概率小。

以随机信道 $H$、设计变量 $x$ 和速率 $R(x,H)$ 为例，三种要求是
\begin{align}
 &\max_{x\in C}\ \E[R(x,H)], &&\text{平均速率最大化},\\
 &\Prob\{R(x,H)<r_0\}\le\epsilon, &&\text{中断概率约束},\\
 &R(x,H)\ge r_0\quad\text{对所有 }H\in\mathcal H, &&\text{集合内最坏情形保证}.
\end{align}
其中 $C$ 是确定的可行集，$r_0$ 是目标速率，$\epsilon\in(0,1)$ 是允许的中断概率，$\mathcal H$ 是误差集合。若 $x$ 在信道实现前选定，它是一个固定设计；若发射机每时隙先观测 $H$ 再决定功率，决策对象变成策略 $x(H)$。后者还需要说明能获得哪些观测以及是否有平均功率约束，不能把两种信息结构混在同一个模型中。

\section{样本平均：用有限数据近似期望}
当期望没有方便的解析表达时，样本平均近似（SAA）把随机目标变为确定优化。对成本函数 $f(x,H)$，从建模分布独立同分布地抽取 $H_1,\ldots,H_N$，求解
\begin{equation}
 \min_{x\in C} F_N(x),\qquad
 F_N(x)=\frac1N\sum_{i=1}^N f(x,H_i).
 \label{eq:saa}
\end{equation}
若每个 $f(\cdot,H_i)$ 都凸，平均后仍然凸，因此可以直接调用前面的凸优化方法。若单个速率最大化本来非凸，增加样本不会消除该结构。

大数定律先保证每个固定 $x$ 的样本平均逼近期望。要把这个结论用于由样本选出的最优解，还需要控制整个可行集上的误差。例如，$C$ 紧、$f(\cdot,H)$ 几乎处处连续，且存在可积的随机上界支配 $\sup_{x\in C}\abs{f(x,H)}$ 时，可在通常的可测性条件下建立一致收敛，进而得到最优值收敛和最优解聚点的有效性。实际所需样本数还受变量维数、分布尾部和目标曲率影响，没有一个对所有通信模型都适用的固定数目\cite{shapiro2009}。

\paragraph{为什么训练样本不能代替验证样本。}
求解器已利用训练样本选择 $x_N$，在同一数据上检查速率会偏向这个选择。应固定设计后，另取 $M$ 个独立验证信道，报告平均速率和中断比例。若真实中断概率为 $p$，独立验证中的中断数为 $V_M$，记 $\widehat p=V_M/M$。Hoeffding 不等式给出一个简单上界：对置信参数 $\delta\in(0,1)$，
\begin{equation}
 p\le\min\!\left\{1,\widehat p+
 \sqrt{\frac{\ln(1/\delta)}{2M}}\right\}
 \quad\text{以至少 }1-\delta\text{ 的概率成立}.
 \label{eq:validation-bound}
\end{equation}
这里的概率针对验证数据的随机抽取，$p$ 则是已选设计在信道分布下的失效概率。$\delta$ 与系统允许的中断概率 $\epsilon$ 含义不同。

若验证中一次中断都没有，可以使用更直接的单侧二项分布上界 $p\le1-\delta^{1/M}$，置信水平为 $1-\delta$。零次失效仍对应非零的概率上界。如果反复查看验证结果并据此调整设计，原先“设计与验证样本独立”的条件就被破坏，应保留新的最终测试集或采用适合该流程的统计方法。时间相关的实测信道也需要相应的相关性分析，不能直接套用独立样本公式。

\section{机会约束：一个可精确变成二阶锥的例子}
线性高斯不确定性下，单个机会约束可以精确写成二阶锥约束。设实随机向量 $a\sim\mathcal N(\bar a,\Sigma)$，其中 $\bar a$ 为均值，$\Sigma\succeq0$ 为协方差；设计向量 $x\in\R^n$ 在 $a$ 实现前确定。对 $0<\epsilon<1/2$，有\cite{ben2009}
\begin{equation}
 \Prob\{a^\T x\le b\}\ge1-\epsilon
 \quad\Longleftrightarrow\quad
 \bar a^\T x+\Phi^{-1}(1-\epsilon)\norm{\Sigma^{1/2}x}_2\le b,
 \label{eq:gaussian-chance}
\end{equation}
其中 $b$ 为确定阈值，$\Phi$ 是标准正态分布函数，$\Sigma^{1/2}$ 可取对称半正定平方根。概率越严格，正分位数越大，设计就需保留更大的安全余量。

证明来自一个标量事实：$a^\T x$ 的均值为 $\bar a^\T x$，方差为 $x^\T\Sigma x$。方差非零时标准化后应用正态分位数，即得右式；方差为零时它退化成确定约束。因为分位数非负，右式是凸的。这个推导适用于线性高斯量；随机二次 SINR 或非高斯衰落下的中断约束需要另行分析，不能直接替换一个协方差矩阵就使用同一公式。

\paragraph{逐用户可靠与全系统可靠。}
若有 $K$ 个用户，分别要求事件 $E_k$ 成立的概率至少 $1-\epsilon$，并不能得到所有用户同时满足的概率也至少 $1-\epsilon$。要保证
$\Prob(\bigcap_k E_k)\ge1-\epsilon$，一个容易实现的充分条件是给用户分配 $\epsilon_k>0$，使
\begin{equation}
 \sum_{k=1}^K\epsilon_k\le\epsilon,\qquad
 \Prob(E_k^c)\le\epsilon_k.
\end{equation}
由并集界，$\Prob(\bigcup_kE_k^c)\le\sum_k\epsilon_k\le\epsilon$。它不要求用户失效独立，但可能保守。固定 $\epsilon_k$ 时，每个线性高斯约束可用式~\eqref{eq:gaussian-chance}；把概率预算也作为变量后，整体凸性需重新检查。

\section{随机梯度：每次只处理一小批信道}
若信道样本不断到来，或全部样本的梯度计算过重，可以直接使用随机梯度。设 $F(x)=\E[f(x,H)]$ 为凸成本，$C$ 为闭凸可行集；在第 $t$ 轮获得随机方向 $g_t$，更新
\begin{equation}
 x^{t+1}=\Pi_C(x^t-\eta_t g_t).
 \label{eq:sgd}
\end{equation}
令 $\mathcal F_t$ 表示本轮抽样前的全部历史。常用条件是
$\E[g_t\mid\mathcal F_t]\in\partial F(x^t)$，即条件均值是一个次梯度，以及
$\E[\norm{g_t}_2^2\mid\mathcal F_t]\le G^2$。对光滑目标，在可交换微分与期望且抽样正确时，$\nabla_x f(x^t,H_t)$ 可以满足无偏条件。

下面的误差界解释了为什么应平均迭代。固定步长 $\eta$，若最优解 $x^\star$ 存在，$\norm{x^0-x^\star}_2\le D$，令 $\bar x_T=T^{-1}\sum_{t=0}^{T-1}x^t$，则
\begin{equation}
 \E[F(\bar x_T)]-F(x^\star)
 \le\frac{D^2}{2\eta T}+\frac{\eta G^2}{2}.
 \label{eq:sgd-bound}
\end{equation}
由投影的非扩张性展开 $\norm{x^{t+1}-x^\star}_2^2$，再用凸性和无偏条件求条件期望；对 $t$ 求和后，距离项相消，最后对平均点应用凸性即可得到该式。若已知计算轮数 $T$，取 $\eta=D/(G\sqrt T)$ 得到 $DG/\sqrt T$ 的上界。固定步长长期使用时，式中的第二项不会消失。

若希望持续运行并让随机误差逐渐减弱，常用 $\sum_t\eta_t=\infty$、$\sum_t\eta_t^2<\infty$ 的步长条件，例如 $\eta_t=c/(t+1)^a$、$c>0$、$1/2<a\le1$；具体收敛陈述还应核查目标和噪声的有界性条件。增加小批量通常降低方向方差，但每轮计算量更大。对非凸期望目标，合适的一阶条件度量可以趋于零，不能沿用式~\eqref{eq:sgd-bound} 的全局最优误差结论。

\section{队列把长期流量约束变成逐时隙决策}
固定平均速率约束没有描述突发业务等待多久。队列模型让调度器看到尚未服务的数据，并据此在速率与能耗之间取舍。令 $Q_k(t)$ 为时隙 $t$ 开始时用户 $k$ 的队列长度，$A_k(t)$ 为该时隙末到达的分组数，$\mu_k(t)$ 为调度提供的服务量，均以分组为单位，则
\begin{equation}
 Q_k(t+1)=\max\{Q_k(t)-\mu_k(t),0\}+A_k(t).
 \label{eq:queue-update}
\end{equation}
空队列无法消耗的服务被浪费。假设各时隙的外生信道与到达过程独立同分布，$0\le A_k(t)\le A_k^{\max}$，$0\le\mu_k(t)\le\mu_k^{\max}$，功率成本 $P(t)\in[0,P^{\max}]$，并存在满足瞬时约束的平稳随机策略，使每个用户的平均服务至少超过平均到达率一个固定量 $\varepsilon>0$。这个严格余量条件表示业务到达率位于可支持区域内部。

定义二次 Lyapunov 函数 $L(Q)=\tfrac12\sum_kQ_k^2$，它衡量总积压。对式~\eqref{eq:queue-update} 平方并放缩，可得条件漂移
\begin{equation}
 \Delta(t)\le B+\sum_k Q_k(t)
 \E[A_k(t)-\mu_k(t)\mid Q(t)],\qquad
 B=\frac12\sum_k[(A_k^{\max})^2+(\mu_k^{\max})^2].
\end{equation}
其中 $\Delta(t)=\E[L(Q(t+1))-L(Q(t))\mid Q(t)]$。在漂移上加 $V\E[P(t)\mid Q(t)]$，$V>0$ 为权衡参数。对给定的队列与当前信道，删去动作无关项后，逐时隙只需求解\cite{neely2010}
\begin{equation}
 \max_{\text{当前可行动作}}\quad
 \sum_k Q_k(t)\mu_k(t)-V P(t).
 \label{eq:dpp}
\end{equation}
积压大的用户获得更高的速率权重；$V$ 越大，调度越重视节省功率。它不要求先知道到达率，但仍需知道或计算当前动作产生的服务与功率。

若每轮精确最小化上述漂移上界，在有界性、独立同分布外生状态和严格余量等标准条件下，长期平均功率满足
\begin{equation}
 \limsup_{T\to\infty}\frac1T\sum_{t=0}^{T-1}\E[P(t)]
 \le P^\star+\frac BV,
 \label{eq:dpp-cost}
\end{equation}
其中 $P^\star$ 是稳定队列所需的最优平均功率。用有严格余量的比较策略还可得到平均总积压不超过 $(B+VP^{\max})/\varepsilon$ 的界，因此呈现 $O(1/V)$ 的功率差与 $O(V)$ 的积压权衡。初始队列有限时，它对长期平均的附加项随 $T$ 消失。这里的保证来自漂移求和与比较策略，并非单次优化的全局最优性替代物。

通信干扰可能使式~\eqref{eq:dpp} 本身非凸。若每轮只得到驻点，不能直接声称满足式~\eqref{eq:dpp-cost}；若能证明每轮的加性求解误差至多为固定 $C$，则通常把相应界中的 $B$ 改为 $B+C$。队列稳定也不等于逐分组时延约束。在满足平稳性、稳定性和有限均值等 Little 定律条件时，平均队列长度除以有效平均到达率才对应平均时延；时延尾概率还需另外分析。

\section{已知状态转移时：马尔可夫决策模型}
若当前动作会改变未来状态，例如重传次数、电池电量或缓存占用，逐时隙模型还可以组织成马尔可夫决策过程（MDP）。设状态集合 $\mathcal S$ 和各状态的动作集合 $\mathcal A(s)$ 均有限，$P(s'\mid s,a)$ 是采取动作 $a$ 后转移至 $s'$ 的已知概率；$r(s,a)$ 为本轮收益，$c(s,a)$ 为功率等成本。有限缓存和离散电池可构成有限状态模型；把本来无界的队列截断则引入了新的模型近似。

无平均成本约束时，平均收益动态规划的基本方程为
\begin{equation}
 g+h(s)=\max_{a\in\mathcal A(s)}
 \left\{r(s,a)+\sum_{s'\in\mathcal S}P(s'\mid s,a)h(s')\right\}.
\end{equation}
在适当的单链条件下，$g$ 是最优长期平均收益，$h(s)$ 是状态相对价值：同样的即时速率，若能留下更有利的电池或重传状态，就有更高的总价值。这是动态规划与只比较当前速率的区别\cite{puterman1994}。

平均约束也可通过稳态占用测度变成线性规划。令 $z(s,a)$ 表示长期处于状态 $s$ 且选择动作 $a$ 的时间比例。若平稳策略的长期平均不依赖初始状态，例如模型中每个平稳策略均满足单链性质，则已知模型下的平均收益问题可写为
\begin{align}
 \max_{z\ge0}\quad&\sum_{s,a}z(s,a)r(s,a),\\
 \text{s.t.}\quad&\sum_a z(s,a)=
       \sum_{s',a'}P(s\mid s',a')z(s',a'),\quad s\in\mathcal S,\\
 &\sum_{s,a}z(s,a)=1,\qquad
 \sum_{s,a}z(s,a)c(s,a)\le\bar c.
 \label{eq:mdp-occupation}
\end{align}
这里 $\bar c$ 为允许的平均成本；每个求和中的动作均属于相应状态的可行动作集合。流平衡约束使进入状态 $s$ 的长期概率质量等于离开它的质量。由最优 $z$ 恢复随机策略为
$\pi(a\mid s)=z(s,a)/\sum_bz(s,b)$，分母为零的状态可指定任意可行动作，并按模型的单链条件处理其暂态行为。若存在多个互不连通的常返类，初始可达性会影响可实现的占用测度，需要额外约束或分解处理。

占用测度 LP 依赖已知转移模型，且变量数随状态--动作对数增长；多个用户的缓存和电池会使状态数迅速膨胀。Lyapunov 方法利用队列结构构造可分析的在线近似，通常无需完整转移表，但其功率--积压保证与精确有限状态 MDP 的最优性属于不同结论。

\section{设计时间尺度与检验指标相配}
天线配置、用户关联和长时功率预算可以较慢更新；瞬时波束和调度需要跟随信道及队列变化。两时间尺度设计用慢变量规定可用资源，再让快变量在每个时隙分配这些资源。若使用慢、快两组随机迭代，通常需让慢步长与快步长之比趋于零，并核查快子系统的稳定性，才能应用相应的两时间尺度收敛结论。

评价此类方法时，应分别报告离线样本目标、独立样本上的可靠性、时间平均功率以及队列指标。这样能识别差距来自模型分布、有限样本、数值求解还是业务动态，也能判断下一步应增加样本、改进优化算法，还是重新选择可靠性目标。
